%======================================================================
%  Vorlage
%======================================================================
%  $Id$
%  Matthias Kupfer

%======================================================================
%  Documentclass for Koma Script >= 3.0
%======================================================================
% Anleitung unter ftp://ftp.dante.de/pub/tex/macros/latex/contrib/koma-script/scrguide.pdf
\documentclass[
% draft
 fontsize=11pt,  % KOMA default
 a4paper,        % DIN A4
 twoside,        % Zweiseitig
 german,         %
  headsepline=true,  % Linie unter der Kopfzeile  
  footsepline=false, % Linie über Fussnote ja/nein
  automark,      % Kolumnentitel lebendig
  headings=normal, % Überschriften normal setzen
  appendixprefix,  % Anhang
  open=any,      % 
  cleardoublepage=plain, %
  abstract=true, %
  index=totoc,   % Index soll im Inhaltsverzeichnis auftauchen
  listof=totoc,  %
  bibliography=totoc,  %
  BCOR8mm,   %%?? Bindungskorrektur: 
%             % BCOR<Breite des Bindungsverlustes> 
]{scrreprt}
%======================================================================
%  Bearbeitung sowohl mit LaTeX als auch mit pdfLaTeX ermoeglichen
%======================================================================
%\newif\ifpdf 
  %\ifx\pdfoutput\undefined 
  %\pdffalse       % we are not running pdflatex 
%\else 
  %\pdfoutput=1       % we are running pdflatex 
  %\pdftrue 
%\fi


%======================================================================
%  Verwendete Pakete
%======================================================================

\usepackage[ngerman]{babel}    % Sprachen

\usepackage[utf8]{inputenc}   % Eingabe von utf8! Bitte beim Speichern der Dokumente beachten!

% Nur geometry Package benutzen um Seitenränder zu ändern
\usepackage[]{geometry}  % http://mirror.informatik.uni-mannheim.de/pub/mirrors/tex-archive/macros/latex/contrib/geometry/geometry.pdf

%\usepackage[''Optionen'']{acronym}

\usepackage{
  acronym,    % Verwaltung von Abkuerzungen
  bibgerm,    % Deutsche Bibliographie 
  calc,      % Erweiterung der arithmetischen Funktionen in 
  %amsmath, % stop errors while using \AA in math mode, but simply doesn't work
          % LaTeX. wird verwendet um Titelseite zu zentrieren
          % Help: http://tug.ctan.org/tex-archive/macros/latex/required/tools/calc.pdf
  color,      % im Laufenden Text einfach mit \color{Farbe) zwischen den 
          % Farben umschalten, wobei Farbe einfach 
          % durch z.B. red, blue, black etc. ersetzt wird
          % \textcolor{farbe){Text)
%  epigraph,    % Zitat am Kapitelanfang
%  fancyhdr,    % Kopf- und Fußzeilen von Dokumenten frei 
        % gestalten
%  fancybox,    % shadowbox, doublebox, ovalbox, Ovalbox 
%  fancyvrb,    % verbatim Erweiterung:
%  float,      % Positionierung von Gleitobjekten genau an der Stelle, wo man
          % 'figure'- oder 'table'-Umgebung die 
          % Positionierung [H] gesetzt werden
          % Help: /usr/share/texmf/doc/latex/float/float.dvi
%  glosstex,    % Glossar und Abkürzungsverzeichnis
%  scrdate,    % \todaysname 
%  scrtime,    % \thistime
  scrpage2,    % Kopf- und Fußzeilen flexibel gestalten
          % Help: siehe scrguide.pdf
  tabularx,    % Blocksatzspalten
          % Help: http://www.cs.brown.edu/system/software/latex/doc/tabularx.pdf http://de.wikibooks.org/wiki/LaTeX-W%C3%B6rterbuch:_tabularx
%  multicol,    % mehrspaltige Zeilen
%  varioref,    % einheitliche Verweise
%  endnotes,    % Fussnoten -> Endnoten
%  rotating,    % sidewaystable und sidewaysfigure
%  natbib,      % Bibliographie ohne Klammer etc.
 capt-of,
 longtable,
 amsmath,
 pdfpages,
 graphicx,
 %url,   % supposed to show links in literatur
}
\usepackage[rightcaption]{sidecap}


%\usepackage{scrhack} % removes a warning, comes with koma script > 3.0

\usepackage{listings}  % Quellcode schön einbinden
            % ftp://ftp.ctan.org/tex-archive/macros/latex/contrib/listings/listings.pdf
\lstloadlanguages{[LaTeX]TeX}

%\usepackage{setspace}  % Durchschuß, Zeilenabstand
%\doublespace    % doppelzeilig oder
%\onehalfspacing  % anderthalbzeilig

% Für schöne Darstellung von Algorithmen
%\usepackage[german, algoruled, algochapter]{algorithm2e} % http://www.ctan.org/tex-archive/macros/latex/contrib/algorithm2e/algorithm2e.pdf
%\usepackage{algorithmicx} % http://tug.ctan.org/tex-archive/macros/latex/contrib/algorithmicx/algorithmicx.pdf

%======================================================================
%  Bilder, Links
%======================================================================

\usepackage{graphicx}  % http://www.ctan.org/tex-archive/macros/latex/required/graphics/grfguide.pdf
\graphicspath{{img/},{img-ergebnisse/},{img-anhang/},{img-eyecatcher/}}  % Angabe der Pfade, wo die Grafiken liegen; 
            % mehrere Pfade sind möglich

%\usepackage{thumbpdf}  % http://www.ctan.org/tex-archive/support/thumbpdf/readme.txt
%\usepackage[plainpages=false,hidelinks,pdfpagelabels]{hyperref} % http://www.ctan.org/tex-archive/macros/latex/contrib/hyperref/doc/manual.pdf
\usepackage[plainpages=false,pdfpagelabels]{hyperref} % http://www.ctan.org/tex-archive/macros/latex/contrib/hyperref/doc/manual.pdf
\usepackage{amssymb}
\usepackage{subfigure}
\usepackage{scrhack}

%======================================================================
%  user.sty
%======================================================================

\usepackage{user}  % Makros in user.sty 
          % \epsinc{bild.eps}{scale=1}{Bildunterschrift}
          % \missing{Beweis fehlt noch}
          % \comment{Ein Kommentar}

%======================================================================
%  Einstellungen
%======================================================================

\pagestyle{scrheadings}  % Standart  Kopf- und Fußzeile
\setkomafont{pageheadfoot}{\small\scshape} % for new Koma Script

% ---------------------------------------------------------------------
%\setcounter{secnumdepth}{2}
%\setcounter{chapter}{-1}
\setcounter{tocdepth}{3}

% ---------------------------------------------------------------------
%\sloppy    % weniger Worttrennungen, größere Wortabstände
\fussy      % viele Worttrennungen, "schönere" Wortabstände

% ---------------------------------------------------------------------
%\flushbottom      % Ausrichtung der Seitenenden jeweils auf 
        % gleicher Höhe

% ---------------------------------------------------------------------
%\sloppypar    % Das hier relaxt die Einstellungen zum Wortabstand 
      % extrem. Damit ragen keine Worte über den rechten 
      % Zeilenabstand hinaus. Dafür muß stärker auf
      % Wortabstand geachtet werden, der kann dann ziemlich 
      % groß werden. Man erhält aber keine Meldung mehr 
      % über underfull boxes.

%Hiermit kann man das gleiche mit weniger Holzhammer erreichen:
%\setlength{\tolerance}{2000}           % Strafpunkt für Zeilenumbruch
%\setlength{\emergencystretch}{3pt}     % Soweit dürfen einzelne Worte mehr 
          % auseinandergezogen werden
%\setlength{\hfuzz}{1pt}                % Macht den rechten Rand um bis zu 1pt 
          % flatterig.

% ---------------------------------------------------------------------
%Vermeiden einzelner Zeilen am Ende einer Seite oder oben auf einer neuen Seite
\clubpenalty10000
\widowpenalty10000

\hyphenation{konnten hexa-gonalen spin-entartung damit wobei function}
%======================================================================
%  includeonly
%======================================================================

\includeonly{    % Gibt an, welche Dateien der include-Befehl 
          % tatsächlich einfuegen darf.
  metadaten      % Variablen setzen
  ,titel      % Titelseite, Zusammenfassung und Inhaltsverzeichnis
  ,einleitung
  ,grundlagen_cnt
  ,grundlagen_transistor
  ,grundlagen_cntfet
  ,grundlagen_simulation
  ,grundlagen_dft
  ,grundlagen_eht
  ,grundlagen_negf
  ,grundlagen_atk
  ,grundlagen_nds
  ,grundlagen_bte
  ,ergebnisse
  ,auswertung
  % Alle per include einzulesenden Dateien müssen hier angegeben sein!
}

%======================================================================
% * *   T E X T    S T A R T S    H E R E   * * * * * * * * * * * * * *
%======================================================================

\begin{document}
%======================================================================
% Metadaten
%======================================================================
% $Id$
% Original template: Matthias Kupfer
%======================================================================

\newcommand{\dcsubject}{Masterarbeit}
\newcommand{\dctitle}{Atomistische Modellierung und Simulation des Filmwachstums bei Gasphasenabscheidungen}
%% \newcommand{\dctitle}{Modellierung und Simulation von Schichtabscheidungen aus der Gasphase}
\newcommand{\dcsubtitle}{} % Untertitel, falls erforderlich

\newcommand{\dcauthortitle}{B.Sc.}
\newcommand{\dcauthorlastname}{Lorenz}
\newcommand{\dcauthorfirstname}{Erik E.}
\newcommand{\dcauthoremail}{erik.e.lorenz@gmail.com}
\newcommand{\dcdate}{\today}

\newcommand{\dcplace}{Chemnitz} % Ort, kann an der TU meist so bleiben
\newcommand{\dcuni}{Technische Universität \dcplace}
\newcommand{\dcdepart}{Fakultät für Naturwissenschaften} % Fakultätsangabe
\newcommand{\dcinstitute}{Institut für Physik}

\newcommand{\dcinstituteext}{Fraunhofer-Institut für Elektronische Nanosysteme}
\newcommand{\dcdepartext}{Abteilung Back-End of Line} % Fakultätsangabe

\newcommand{\dcexaminerfirst}{Prof. Dr. Stefan E. Schulz}
\newcommand{\dcdepartfirst}{Fakultät für Elektrotechnik und Informationstechnik} % Fakultätsangabe
\newcommand{\dcproffirst}{Technologien der Nanoelektronik} % Angabe der Professur

\newcommand{\dcexaminersecond}{Prof. Dr. Karl Heinz Hoffmann}
\newcommand{\dcdepartsecond}{Fakultät für Physik} % Fakultätsangabe
\newcommand{\dcprofsecond}{Theoretische Physik, insbesondere Computerphysik} % Angabe der Professur
\newcommand{\dcinstitutesecond}{\dcinstitute}

\newcommand{\dcadvisor}{Dr. Jörg Schuster}
\newcommand{\dcinstituteadvisor}{\dcinstituteext}
\newcommand{\dcdepartadvisor}{\dcdepartext}
\newcommand{\dcgroupadvisor}{Simulationsgruppe}

\newcommand{\dckeywords}{Dünne Schichten, chemische Gasphasenabscheidung (CVD), physikalische Gasphasenabscheidung (PVD), Atomlagenabscheidung (ALD), Multiskalenmodellierung, Kinetic Monte Carlo (KMC), Molekulardynamik (MD), ReaxFF}

%%======================================================================
% Einstellungen des Hyperref-Paketes
\hypersetup{%
 pdftitle = {\dctitle}, %
 pdfsubject = {\dcsubject, \dcdate}, %
 pdfauthor = {\dcauthorfirstname~\dcauthorlastname, \dcauthoremail}, %
 pdfkeywords = {\dckeywords}, %
 pdfcreator = {pdfTeX with Hyperref and Thumbpdf}, %
 pdfproducer = {LaTeX, hyperref, thumbpdf}, %
 % weitere pdf-Einstellungen in hyperref.cfg
}
%%======================================================================
 % Variablen, \hypersetup, etc.

%======================================================================
%  Titelseite (Zusammenfassung und Inhaltsverzeichnis)
%======================================================================

%======================================================================
%  Titelseite
%======================================================================

% Wir setzten hier die Seitennummerierung auf Großrömisch, auch wenn diese
% nicht auf dem "Papier" erscheint. Es dient nur der internen Zählung der
% Seiten im hyperref-Paket, welches für die Seitennummerierung in den pdfs
% zuständig ist. Ohne diesen Trick kommen ein paar sinnlose Warnungen und
% die Seiten-Navigation im Adobe Reader kann durcheinander geraten.
\pagenumbering{Roman}


%%======================================================================
%% Schmutztitel
%%======================================================================
%\extratitle{
%  \usekomafont{disposition}\mdseries
%  \begin{center}
%    \Huge \dcsubject\\[1.5ex]
%    \hrule
%    \vspace*{\fill}
%    \includegraphics{TUC_deutsch_einzeile_CMYK}
%  \end{center}
%}

%%======================================================================
%% Titelkopf
%%======================================================================
\titlehead{
  \vspace*{-1.5cm}
  \begin{center}
    \raisebox{-1ex}{\includegraphics[scale=1.1]{img/logo_tu_tailored}}\\
    \hrulefill \\[1em]
    {\Large\dcdepart}\\[0.2em]
    \dcinstitute\\[0.5em]
%    {\Large\dcdepartsecond}\\[0.5em]
%    \dcprofsecond
  \end{center}
  \vspace*{0.1cm}
  \begin{center}
    \raisebox{-1ex}{\includegraphics[scale=1.2]{img/logo_fhg_tailored}}\\
    \hrulefill \\[1em]
    {\Large\dcinstituteext}\\[0.2em]
    \dcdepartext\\[0.5em]
  \end{center}
  \vspace*{0.5cm}
}



%%======================================================================
%% Subjekt
%%======================================================================
\subject{\huge\textnormal{\textsc{\dcsubject}}}


%%======================================================================
%% Titel
%%======================================================================
\title{\huge
  \dctitle
  \\
  \dcsubtitle
  \vspace*{-0.4em}
}

%%======================================================================
%% Autor des Dokumentes
%%======================================================================
\author{\dcauthortitle~\dcauthorfirstname~\dcauthorlastname}

%%======================================================================
%% Ort, Datum
%%======================================================================
\date{\dcplace, den \dcdate
}


%%======================================================================
%% Publishers
%%======================================================================
\publishers{
    \begin{tabbing}
      \hspace*{10.1em}\=\kill
      \hspace*{4.8em}{Gutachter:}\>\dcexaminerfirst\\[-0.2em]
                  \>{\small\dcproffirst}\\[0.2em]
                  \>\dcexaminersecond\\[-0.2em]
                  \>{\small\dcprofsecond}\\
\\
      \hspace*{4.8em}{Betreuer:}\>\dcadvisor\\[-0.2em]
                  \>{\small\dcinstituteadvisor}
    \end{tabbing}
  \vspace*{-7em}
}

%%======================================================================
%% bibliografische Angaben
%%======================================================================
\lowertitleback{
\textbf{\dcauthorlastname, \dcauthorfirstname}\\
\textit{\dctitle}\\
%\dcsubject,~\dcdepart\\
\dcsubject\\
\dcuni,~\ifcase\month\or
  Januar\or Februar\or März\or April\or Mai\or Juni\or
    Juli\or August\or September\or Oktober\or November\or Dezember\fi
    ~\number\year\\
Stichworte: \dckeywords
}

%%======================================================================
%% maketitle
%%======================================================================

\maketitle

%%======================================================================
%% Danksagung
%%======================================================================
%\thispagestyle{empty}
%\null\vfil
%\begin{center}
%\usekomafont{disposition}\textbf{Danksagung}
%\vspace{-.5em}\vspace{\parsep}
%%
%% Hier steht der Text für die Danksagung
%%
%\end{center}
%\par\vfil\null
%\cleardoubleemptypage

%%======================================================================
%%      Kurzfassung / Abstract
%%======================================================================
%\def\abstractname{Abstract} % Wenn der Text "Zusammenfassung" erscheinen
                             % soll, dann muß dies auskommentiert werden
 % Titelseite, Zusammenfassung 
        % Inhaltsverzeichnis
        % Abkuerzungsverzeichnis
	% Literaturverzeichnis
        % -> danach Seitennummerierung bei 1

\pagenumbering{arabic} % Fäng erneut bei 1 an.
\cleardoublepage

%======================================================================
%  Abbildungsverzeichnis
%======================================================================
%\cleardoublepage

%======================================================================
%  Inhalt per include
%======================================================================

\newpage
\chapter{Einleitung}
\label{cha:einleitung}

\chapter{Grundlagen}
\label{cha:grundlagen}

\section{Kohlenstoffnanoröhrchen}
\label{cha:kohlenstoffnanoroehrchen}

\section{Allgemeiner Überblick über Transistoren}
\label{cha:allgemeinerueberblickuebertransistoren}

\section{Transistoren auf Basis von CNTs}
\label{cha:transistorenaufbasisvoncnts}

\chapter{Simulationsmethoden}
\label{cha:simulationsmethoden}

\section{Über die Notwendigkeit von Simulationen}
\label{ueberdienotwendigkeitvonsimulationen}

Wie bereits in Kapitel \ref{cha:einleitung} angerissen wurde, existieren zahlreiche Anwendungsszenarien für Simulationen im Zusammenhang mit CNTFETs.\\
Zum einen sind sie unerlässlich bei der \textit{Einschätzung neuer Technologien}. Insbesondere aufgrund hoher Kosten, die beim Wechsel der aktuell genutzten Technologie stets anfallen, sind vorangehende Simulationen von größter Bedeutung. Nicht nur die Abschätzung der Leistung der potenziellen Technologie ist hier wichtig, sondern auch die Frage nach Zuverlässigkeit, Lebensdauer als auch der anfallenden Kosten. Nur wenn die neue Technologie diesen und einigen weiteren Ansprüchen genügt, ist ein Wechsel aus wirtschaftlichen Gründen sinnvoll.\\
Auch die \textit{Optimierung der Transistoren} stellt ein wichtiges Anwendungsfeld für Simulationen dar.
Hier gilt es die optimale Kombination zahlreicher möglicher Parameter zu finden.
Neben der Wahl der verwendeten Materialien gehören hierzu auch geometrische Eigenschaften.
Die Ergebnisse können schließlich als Ausgangspunkt für experimentelle Arbeiten als auch industrielle Umsetzungen dienen.\\
Schließlich ist dank Simulationen auch ein \textit{tieferer Einblick in die Physik} der Transistoren möglich, was für das Verstehen experimenteller Ergebnisse notwendig ist.
Dies ist insbesondere aufgrund der zunehmenden Miniaturisierung der Transistoren von Bedeutung, da aufgrund der kleinen Abmessungen atomare sowie quantenmechanische Effekte berücksichtigt werden müssen.
Nur dank eines fundierten theoretischen Verständnisses ist eine korrekte Beschreibung der Transistoreigenschaften sowie des Transistorverhaltens möglich.

\section{Multiskalenmodellierung}
\label{multiskalenmodellierung}

Für die theoretische Beschreibung und Simulation von Transistoren sind zahlreiche Ansätze denkbar.\\
Aufgrund der sehr kleinen Abmessungen der heutigen und insbesondere zukünftiger Transistoren bieten sich \textit{atomistische Simulationen} an.
Hierbei stellt die Position der einzelnen Atome den Ausgangspunkt der Rechnungen dar.
In Abhängigkeit von der verwendeten Methode sind dabei keine -- wie im Falle von ab initio Methoden -- oder nur eine geringe Anzahl externer Parameter notwendig.
Damit verbunden ist jedoch auch ein relativ hoher Rechenaufwand, was diese Methoden meist nur für die Untersuchung kleinerer Systeme praktikabel macht.\\
Eine Alternative stellen sogenannte \textit{numerische Gerätesimulationen} dar.
Anstelle der exakten Atompositionen kommen hier effektive Modelle zur Beschreibung der Systemeigenschaften zum Einsatz.
Hierfür sind jedoch externe Parameter notwendig, welche beispielsweise von atomistischen Rechnungen oder experimentellen Untersuchen extrahiert werden müssen.
Gleichzeitig können unbekannte Parameter jedoch auch als Fit-Parameter genutzt werden, um so Übereinstimmung mit experimentellen Resultaten zu erreichen.
Aussagen für das Experiment können so schnell erzielt werden, eine physikalisch korrekte Beschreibung des Systems ist so jedoch nur begrenzt möglich.
Numerische Gerätesimulationen erlauben jedoch auch eine weitaus schnellere Berechnung der Systeme.
Auch können sie genutzt werden, um Gesetzmäßigkeiten abzuleiten. Bei atomistischen Verfahren ist dies vergleichsweise schwierig, da \todo{Warum?}.\\
Eine Beschreibung basierend auf Gesetzmäßigkeiten und Formeln bildet die Grundlage für sogenannte \textit{Kompaktmodelle}.
Diese können auf physikalischen Überlegungen und Gesetzen oder aber auch nur auf Fit-Prozeduren basieren.\\
Schließlich ermöglichen \textit{Schaltkreissimulation} die Beschreibung ganzer Schaltkreise mit bis zu \todo{??? Literatur?} Transistoren.
Dabei können wie beispielsweise bei \todo{Lit.} die Ergebnisse der Kompaktmodelle direkt als Eingabe verwendet werden.\\ 
In Tabelle \ref{tab:simulationsmethoden} sind einige Implementierungen der vorgestellten Methoden aufgeführt.

\newcolumntype{Y}{>{\centering\arraybackslash}X}
\begin{table}
\caption{Einige Implementationen der unterschiedlichen Simulationsmodelle, welche im Text vorgestellt wurden.}
\begin{tabularx}{\linewidth}{|Y|Y|Y|Y|}
\hline
Atomistische Simulation & Numerische Gerätesimulation & Kompakt\-modell & Schaltkreis\-simulation \\
\hline
Zitat & Zitat & Zitat & Zitat \\
\hline
\end{tabularx}
\label{tab:simulationsmethoden}
\end{table}

\section{Dichtefunktionaltheorie}
\label{cha:dichtefunktionaltheorie}

\subsection{Erweiterte Hückelmethode}
\label{cha:erweitertehueckelmethode}

\section{Nichtgleichgewichts-Greensfunktions-Formalismus}
\label{cha:nichtgleichgewichtsgreensfunktionsformalismus}

\section{Atomistix ToolKit}
\label{cha:atomistixtoolkit}

\section{Kontinuumsbeschreibung}
\label{kontinuumsbeschreibung}

\subsection{Numerische Gerätesimulation}
\label{cha:numerischegeraetesimulation}

\section{Boltzmann-Transport-Gleichungs-Löser}
\label{cha:boltzmanntransportgleichungsloeser}

\chapter{Berechnungen und Ergebnisse}
\label{cha:ergebnisse}

\chapter{Zusammenfassung der Ergebnisse und Ausblick}
\label{cha:auswertung}


%======================================================================
%  Glossar
%======================================================================

%\manualmark
%\addcontentsline{toc}{chapter}{Glossar}
%\markboth{Glossar}{Glossar}
%\def\glossaryname{Glossar}
%\printglosstex(glo)
%\cleardoublepage

%======================================================================
%  Literaturverzeichnis
%======================================================================
\cleardoublepage
%\manualmark
%\markboth{Literaturverzeichnis}{Literaturverzeichnis}
%\bibliographystyle{unsrt} % standart
\bibliographystyle{gerunsrtmod} % titel kursiv, journal normal; ist fein, modifiziert damit die Namen stimmen
%\bibliographystyle{gerunsrtdoi} % titel kursiv, journal normal; ist fein, modifiziert damit die Namen stimmen + doi anzeigen (in progress ...)
%\bibliographystyle{IEEEtrans} % Hinweis von Christian
%\usepackage{url}

\bibliography{literatur} %standart


%\cleardoublepage

%\newpage
%\cleardoublepage

%======================================================================
%  Anhang
%======================================================================

\begin{appendix}
%\include{anhanga}
%\newpage
\include{anhangb}
\include{anhangc}
%\addcontentsline{toc}{chapter}{B. Danksagung}
\thispagestyle{empty}
\chapter{Danksagung}
\newpage
%\chapter{}
\null\vfil
\begin{center}
%\vspace{-.5em}\vspace{\parsep}
%\vspace{0.2cm}
\vspace{3cm}
\end{center}
\par\vfil\null
%\cleardoubleemptypage

%======================================================================
%  Selbstständigkeitserklärung
%======================================================================

%\addcontentsline{toc}{chapter}{C. Selbständigkeitserklärung}
\chapter{Selbstständigkeitserklärung}
%\includepdf[pagecommand={}]{selbststaendigkeitserklaerung.pdf}
%\stepcounter{page}
%\includepdf{selbststaendigkeitserklaerung.pdf}
\cleardoublepage
\includepdf[pagecommand={\thispagestyle{plain}}]{selbststaendigkeitserklaerung.pdf}

\end{appendix}

%\part*{Anhang}

%\include{anhanga}

\end{document}
